% ---------------------------------------------------------------------------
% Định dạng báo cáo. Biên dịch bằng XeLaTeX (cần cho tiếng Việt + font hệ thống).
% ---------------------------------------------------------------------------

\usepackage{amsmath}
\usepackage{fontspec}
\usepackage[a4paper,left=30mm,right=20mm,top=25mm,bottom=22mm,headheight=15pt,
            footskip=12mm]{geometry}
\usepackage{setspace}
\usepackage{microtype}
\usepackage{graphicx}
\usepackage{xcolor}
\usepackage{booktabs}
\usepackage{array}
\usepackage{longtable}
\usepackage{tabularx}
\usepackage{caption}
\usepackage{enumitem}
\usepackage{titlesec}
\usepackage{tocloft}
\usepackage{fancyhdr}
\usepackage{fancyvrb}
\usepackage{tcolorbox}
\tcbuselibrary{breakable,skins}
\usepackage{needspace}
\usepackage[hidelinks]{hyperref}
\usepackage{xurl}

% --- Font ------------------------------------------------------------------
% Cả ba font đều thuộc TeX Gyre / TeX Live nên có sẵn ở mọi bản cài, và quan trọng
% hơn: cả ba đều phủ ĐỦ bảng chữ tiếng Việt dựng sẵn (U+1EA0–U+1EF9).
%   Pagella       — bản TeX Gyre của Palatino, thân bài, dễ đọc khi in dài.
%   Heros         — bản TeX Gyre của Helvetica, tiêu đề/bảng/chú thích; nét sạch,
%                   tương phản rõ với Pagella mà không ồn như font geometric.
%   LM Mono       — mã nguồn. KHÔNG dùng DejaVu Sans Mono: font đó thiếu 21/49 ký
%                   tự tiếng Việt dựng sẵn nên HarfBuzz phải tách thành ký tự nền +
%                   dấu rời, và trong font đơn cách dấu chiếm hẳn một ô -> vỡ chữ
%                   ("nếu" hiện thành "nêú").
\setmainfont{TeX Gyre Pagella}[Ligatures=TeX]
\setsansfont{TeX Gyre Heros}[Ligatures=TeX, Scale=0.92]
\setmonofont{Latin Modern Mono}[Scale=0.95]

% --- Màu -------------------------------------------------------------------
\definecolor{bkblue}{HTML}{1A4F8B}
\definecolor{bkdark}{HTML}{15202B}
\definecolor{bkmuted}{HTML}{5B6B80}
\definecolor{bkrule}{HTML}{C9D4E2}
\definecolor{bksoft}{HTML}{EEF3FA}
\definecolor{bkcode}{HTML}{F4F6F9}

\color{bkdark}
\setstretch{1.12}
\setlength{\parskip}{0.55em}
\setlength{\parindent}{0pt}
\setlength{\emergencystretch}{2em}

% --- Tiêu đề mục -----------------------------------------------------------
\titleformat{\section}[block]
  {\sffamily\bfseries\LARGE\color{bkblue}}{\thesection.}{0.5em}{}
  [\vspace{-0.5em}{\color{bkblue}\rule{\textwidth}{1.2pt}}]
\titlespacing*{\section}{0pt}{0pt}{1.1em}

\titleformat{\subsection}[block]
  {\sffamily\bfseries\large\color{bkdark}}{\thesubsection.}{0.5em}{}
\titlespacing*{\subsection}{0pt}{1.5em}{0.5em}

\titleformat{\subsubsection}[block]
  {\sffamily\bfseries\normalsize\color{bkmuted}}{}{0em}{}
\titlespacing*{\subsubsection}{0pt}{1.2em}{0.35em}

% Mỗi chương bắt đầu một trang mới.
\newcommand{\chapterbreak}{\clearpage}

% --- Mục lục ---------------------------------------------------------------
\renewcommand{\cfttoctitlefont}{\sffamily\bfseries\LARGE\color{bkblue}}
\renewcommand{\cftaftertoctitle}{\par\vspace{-0.3em}{\color{bkblue}\rule{\textwidth}{1.2pt}}}
\renewcommand{\cftsecfont}{\sffamily\bfseries}
\renewcommand{\cftsecpagefont}{\sffamily\bfseries}
\renewcommand{\cftsubsecfont}{\sffamily\small}
\renewcommand{\cftsubsecpagefont}{\sffamily\small\color{bkmuted}}
\renewcommand{\cftsecleader}{\sffamily\cftdotfill{\cftdotsep}}
\renewcommand{\cftsubsecleader}{\cftdotfill{\cftdotsep}}
\setlength{\cftbeforesecskip}{0.7em}
\setlength{\cftbeforesubsecskip}{0.18em}
\setcounter{tocdepth}{2}
\setcounter{secnumdepth}{2}

% --- Đầu trang / chân trang ------------------------------------------------
\pagestyle{fancy}
\fancyhf{}
\renewcommand{\headrulewidth}{0.4pt}
\renewcommand{\headrule}{\hbox to\headwidth{\color{bkrule}\leaders\hrule height \headrulewidth\hfill}}
\fancyhead[L]{\sffamily\scriptsize\color{bkmuted} Hệ thống RAG truy vấn quy định học vụ}
\fancyhead[R]{\sffamily\scriptsize\color{bkmuted}\nouppercase{\leftmark}}
\fancyfoot[C]{\sffamily\small\color{bkmuted}\thepage}
\renewcommand{\sectionmark}[1]{\markboth{#1}{}}

\fancypagestyle{plain}{\fancyhf{}\renewcommand{\headrulewidth}{0pt}%
  \fancyfoot[C]{\sffamily\small\color{bkmuted}\thepage}}

% --- Hình và bảng ----------------------------------------------------------
\captionsetup{font={sf,small},labelfont={sf,bf,color=bkdark},
              labelsep=period,justification=centering,width=0.94\textwidth,
              skip=6pt}
\captionsetup[table]{justification=raggedright,singlelinecheck=false,
                     width=\textwidth,skip=4pt,position=below}
\renewcommand{\figurename}{Hình}
\renewcommand{\tablename}{Bảng}
% Tên do LaTeX tự sinh, đổi sang tiếng Việt.
\renewcommand{\contentsname}{Mục lục}
\renewcommand{\appendixname}{Phụ lục}
\renewcommand{\refname}{Tài liệu tham khảo}
\renewcommand{\arraystretch}{1.25}
\setlength{\tabcolsep}{6pt}

% Bảng theo phong cách booktabs: chỉ ba đường kẻ ngang, không kẻ dọc, không tô nền —
% kiểu trình bày chuẩn của tài liệu học thuật và dễ đọc hơn bảng có khung đầy đủ.
% \tabhead giữ lại như một điểm móc nếu sau này muốn đổi cách làm nổi dòng tiêu đề.
\newcommand{\tabhead}[1]{#1}
\newenvironment{rtable}
  {\par\needspace{5\baselineskip}\begingroup\sffamily\small\centering}
  {\par\endgroup}

% --- Danh sách -------------------------------------------------------------
\setlist[itemize]{leftmargin=1.3em,itemsep=0.25em,topsep=0.35em,parsep=0pt}
\setlist[enumerate]{leftmargin=1.5em,itemsep=0.3em,topsep=0.35em,parsep=0pt}

% --- Mã nguồn --------------------------------------------------------------
\newcommand{\code}[1]{{\ttfamily\small #1}}

\DefineVerbatimEnvironment{codeblock}{Verbatim}
  {fontsize=\small,xleftmargin=10pt,baselinestretch=1.0,
   frame=leftline,framerule=1.6pt,rulecolor=\color{bkblue},framesep=8pt}

% --- Khung ghi chú ---------------------------------------------------------
\newtcolorbox{notebox}[1][]{
  colback=bksoft, colframe=bkblue, boxrule=0pt, leftrule=2.2pt,
  arc=1pt, left=9pt, right=9pt, top=6pt, bottom=6pt,
  breakable, enhanced, before skip=10pt, after skip=10pt,
  fonttitle=\sffamily\bfseries\small\color{bkblue}, coltitle=bkblue,
  attach title to upper={\par\vspace{2pt}}, #1}

% --- Siêu dữ liệu PDF ------------------------------------------------------
\hypersetup{
  pdftitle={Xây dựng hệ thống RAG hỗ trợ truy vấn quy định học vụ},
  pdfauthor={Nguyễn Thái Thành Đạt, Lê Công Tú, Ngô Nhất Toàn},
  pdfsubject={CO5179 — Những xu hướng mới trong lĩnh vực công nghệ thông tin},
  bookmarksnumbered=true, bookmarksopen=true, breaklinks=true}
